%------------------------------------------------------------------------------------
%  SECTION 2 : CITATIONS ET BIBLIOGRAPHIE
%------------------------------------------------------------------------------------
\section{Citations et Gestion Bibliographique}
\label{sec:bibliographie}

Le moteur de gestion bibliographique utilisé est \texttt{biblatex} associé au backend \texttt{biber}.
Le style sélectionné ici est Vancouver, adapté aux publications médicales et scientifiques.

\subsection{Exemples de citations en contexte}
\label{subsec:exemples_citations}

Différents modes de citation sont pris en charge selon la syntaxe de la phrase :
\begin{itemize}
  \item \textbf{Citation intégrée au texte (auteur-date ou sujet) :}
  Comme le démontrent \textcite{nieva2024impact}, l'évolution des recommandations influence directement les diagnostics précoces.
  \item \textbf{Citation globale ou parenthétique :}
  Plusieurs biomarqueurs sanguins prometteurs font l'objet d'évaluations multicentriques récentes~\parencite{burghaus2024}.
  \item \textbf{Citation avec numéro standard :}
  L'échographie transvaginale demeure un examen clé dans la détection des lésions ovariennes~\cite{kanti2024ultrasound}.
  \item \textbf{Sources web et ressources en ligne :}
  Les principes de communication interpersonnelle sont également documentés en ligne~\parencite{demetersante2022}.
\end{itemize}

Les liens hypertextes créés sont directement cliquables et renvoient vers la section bibliographique en fin de document.
De plus, la fonctionnalité \texttt{backref} affiche les pages précises où chaque document a été référencé.

